% !TEX program = tectonic
\documentclass[12pt,a4paper]{article}

% ── Packages ──────────────────────────────────────────────────
\usepackage[utf8]{inputenc}
\usepackage[T1]{fontenc}
\usepackage[margin=2.5cm]{geometry}
\usepackage{hyperref}
\usepackage{xcolor}
\usepackage{listings}
\usepackage{enumitem}
\usepackage{booktabs}
\usepackage{needspace}
\usepackage{parskip}
\usepackage{fancyhdr}
\usepackage{titlesec}

% ── Colors ────────────────────────────────────────────────────
\definecolor{codebg}{RGB}{245,245,245}
\definecolor{codeframe}{RGB}{200,200,200}
\definecolor{comment}{RGB}{128,128,128}
\definecolor{keyword}{RGB}{0,0,180}
\definecolor{string}{RGB}{0,128,0}

% ── Section formatting (16pt centered, 14pt left) ─────────────
\titleformat{\section}[block]{\centering\bfseries\fontsize{16}{19}\selectfont}{}{0pt}{}
\titleformat{\subsection}[hang]{\bfseries\fontsize{14}{17}\selectfont}{}{0pt}{}
\titleformat{\subsubsection}[hang]{\bfseries\fontsize{12}{14}\selectfont}{}{0pt}{}

% ── Page style ────────────────────────────────────────────────
\pagestyle{fancy}
\fancyhf{}
\fancyhead[L]{\small Safe-Concurrency Multi-Sensor Fusion}
\fancyhead[R]{\small Naskah Presentasi}
\fancyfoot[C]{\thepage}
\renewcommand{\headrulewidth}{0.4pt}

% ── Code listing style ────────────────────────────────────────
\lstdefinestyle{data}{
    backgroundcolor=\color{codebg},
    frame=single,
    rulecolor=\color{codeframe},
    basicstyle=\ttfamily\footnotesize,
    breaklines=true,
    breakatwhitespace=false,
    columns=flexible,
    keepspaces=true,
    aboveskip=8pt,
    belowskip=8pt,
}
\lstdefinestyle{prompt}{
    backgroundcolor=\color{codebg},
    frame=single,
    rulecolor=\color{codeframe},
    basicstyle=\ttfamily\footnotesize,
    breaklines=true,
    columns=flexible,
}

% ── Meta ──────────────────────────────────────────────────────
\hypersetup{
    colorlinks=true,
    linkcolor=blue,
    urlcolor=blue,
}

\begin{document}

% ═══════════════════════════════════════════════════════════════
% TITLE PAGE
% ═══════════════════════════════════════════════════════════════
\begin{titlepage}
\centering
\vspace*{3cm}

{\fontsize{18}{22}\selectfont\bfseries Naskah Presentasi \& Bahan Bacaan Mendalam\par}
\vspace{0.8cm}
{\fontsize{14}{17}\selectfont Safe-Concurrency for Multi-Sensor Fusion\par}
{\fontsize{14}{17}\selectfont in Industrial Safety-Critical Systems\par}
\vspace{0.5cm}
{\fontsize{12}{14}\selectfont N-Sensor + Adaptive Lockout + Event-Triggered\par}

\vspace{2cm}
{\large\textbf{Abdurrauf Almutawakkil} (2042241115)\par}
\vspace{0.3cm}
{\large Mata Kuliah: Pemrograman Kontroller (ETS Bagian B)\par}
\vspace{0.3cm}
{\large Teknik Instrumentasi -- ITS\par}

\vspace{2cm}
{\normalsize Untuk dipelajari, bukan sekadar dihafal.\par}
\vspace{0.3cm}
{\normalsize Data asli dari simulasi Proteus, 18 Juni 2026\par}

\vfill
\end{titlepage}

% ═══════════════════════════════════════════════════════════════
\section{BAGIAN 1 --- Narasi Besar (The Big Story)}

\subsection{Apa masalah yang kita pecahkan?}

Di dunia industri, sistem safety-critical seperti kontrol valve pada pabrik kimia,
pembangkit listrik, atau kilang minyak TIDAK BOLEH GAGAL. Satu kesalahan bisa
berakibat fatal: ledakan, kebocoran gas beracun, kerusakan lingkungan.

Saat ini, mayoritas sistem embedded safety-critical ditulis dalam bahasa C/C++.
Masalahnya:
\begin{itemize}[nosep]
    \item C/C++ tidak menjamin memory safety -- programmer bisa membuat bug seperti
          use-after-free, buffer overflow, atau data race.
    \item 186 CVE (vulnerability) terkait memory-safety ditemukan di Rust unsafe code
          saja (Xu et al., 2021) -- apalagi C/C++ yang tidak punya safety guarantee.
    \item Sistem konvensional biasanya menggunakan single-sensor threshold -- satu sensor
          rusak = false positive = shutdown tidak perlu, ATAU false negative = bencana.
\end{itemize}

\subsection{Apa solusi yang kami tawarkan?}

Kami membangun sistem sensor fusion dengan tiga lapisan keamanan:
\begin{enumerate}[nosep]
    \item \textbf{Rust safe-concurrency} -- compiler menjamin tidak ada data race. Bukan
          programmer yang harus hati-hati, tapi compiler yang menolak compile jika
          ada potensi race condition.
    \item \textbf{Voting redundancy} -- fault hanya dideklarasikan jika $\ge 2$ dari 3 sensor
          anomali. Satu sensor rusak masih ditoleransi.
    \item \textbf{Adaptive lockout} -- setelah fault, aktuator dikunci dalam posisi aman.
          Durasi lockout disesuaikan dengan severity: 500ms untuk MINOR (2 sensor),
          2000ms untuk CRITICAL (3 sensor).
\end{enumerate}

% ═══════════════════════════════════════════════════════════════
\section{BAGIAN 2 --- Memahami Setiap Komponen}

\subsection{Rust Safe-Concurrency -- Kenapa Ini Penting?}

\textbf{Analoginya:} Bayangkan tiga orang (sensor) menulis di satu papan tulis
(shared state) secara bersamaan. Di C/C++, programmer harus mengatur siapa
yang pegang spidol -- kalau lupa, dua orang bisa menulis bersamaan dan hasilnya
kacau (data race). Di Rust, compiler BERTINDAK SEPERTI SATPAM yang memastikan
hanya satu orang yang bisa menulis di papan setiap saat.

\textbf{Secara teknis:} Kami menggunakan \texttt{Mutex<RefCell<T>>} yang dibungkus dalam
\texttt{critical\_section}. Begitu critical section dimulai, SEMUA interrupt di
ESP32-S3 dinonaktifkan -- jadi tidak ada yang bisa mengakses data sensor saat
sedang dibaca/ditulis. Begitu critical section selesai, interrupt kembali normal.

\textbf{Ini BUKAN multitasking.} ESP32-S3 dual-core, tapi kami jalan di single-core
bare-metal. Tidak ada RTOS. Concurrency di sini maksudnya: interrupt handler
(timer) bisa fire kapan saja, dan Rust menjamin interrupt handler tidak akan
membaca data yang sedang setengah ditulis.

\needspace{4cm}
\textbf{Pertanyaan yang mungkin muncul:}

\begin{description}[nosep,leftmargin=2em,style=nextline]
    \item[Kenapa tidak pakai RTOS saja?]
    Karena RTOS menambah kompleksitas, overhead, dan potential failure point. Di
    sistem safety-critical, semakin sedikit komponen = semakin sedikit yang bisa gagal.
    Bare-metal artinya KITA yang punya kontrol penuh atas setiap instruksi.

    \item[Bukankah menonaktifkan interrupt itu berbahaya?]
    Ya, makanya critical section harus SANGAT SINGKAT. Kami hanya menonaktifkan
    interrupt selama beberapa mikrodetik -- hanya saat membaca/menulis state sensor.
    Setelah itu interrupt kembali normal. Timer tidak akan miss karena durasinya
    sangat pendek.
\end{description}

\subsection{Voting Redundansi -- Kenapa $\ge 2$ dari 3?}

\textbf{Analoginya:} Tiga orang saksi melihat kejadian. Dua bilang ``terjadi'', satu
bilang ``tidak''. Kita percaya mayoritas. Kalau cuma satu yang bilang ``terjadi'',
mungkin dia salah lihat.

\textbf{Threshold-nya:}
\begin{itemize}[nosep]
    \item Suhu $> 80^{\circ}\mathrm{C}$ $\rightarrow$ anomali (nilai normal: $25^{\circ}\mathrm{C}$)
    \item Tekanan $< 900$ atau $> 1200$ hPa $\rightarrow$ anomali (nilai normal: 1013 hPa)
    \item Vibrasi $> 500$ $\rightarrow$ anomali (nilai normal: 5)
\end{itemize}

\textbf{Kenapa threshold segitu?}
\begin{itemize}[nosep]
    \item $80^{\circ}\mathrm{C}$: di atas suhu operasi normal kebanyakan peralatan industri. Bukan nilai
          random -- didasarkan pada typical industrial temperature rating.
    \item 900--1200 hPa: rentang tekanan atmosfer normal plus margin. Di luar itu
          berarti sensor rusak atau ada kebocoran sistem.
    \item 500: nilai vibrasi signifikan yang mengindikasikan getaran abnormal.
\end{itemize}

\textbf{Pertanyaan yang mungkin muncul:}

\begin{description}[nosep,leftmargin=2em,style=nextline]
    \item[Kenapa tidak 3/3 saja biar lebih aman?]
    Karena kalau 3/3, satu sensor rusak = sistem tidak akan pernah mendeteksi fault
    (selalu maksimal 2/3). Kita justru MAU toleransi 1 sensor gagal -- itulah gunanya
    redundansi.

    \item[Kenapa tidak 1/3? Kan lebih sensitif?]
    Terlalu sensitif = banyak false positive. Satu sensor bisa noise sesaat, dan
    sistem akan shutdown tidak perlu. Di industri, false positive juga berbahaya
    -- bisa menghentikan produksi yang tidak perlu.
\end{description}

\subsection{Adaptive Lockout -- Kenapa 500ms dan 2000ms?}

\textbf{Analoginya:} Seperti airbag di mobil. Begitu mengembang, airbag TIDAK
langsung kempes -- butuh waktu beberapa detik. Kenapa? Karena kalau langsung
kempes, penumpang belum sempat aman, dan tabrakan susulan bisa fatal.

Begitu juga dengan valve industri. Begitu valve menutup karena fault, kita
TIDAK MAU valve langsung membuka lagi hanya karena sensor sudah normal. Kenapa?
\begin{enumerate}[nosep]
    \item Sensor bisa berfluktuasi -- normal $\rightarrow$ anomali $\rightarrow$ normal dalam hitungan milidetik
    \item Proses industri butuh waktu untuk stabil -- buru-buru buka valve bisa
          menyebabkan pressure surge
    \item Valve yang membuka-tutup-membuka-tutup cepat (= valve bounce) bisa merusak
          aktuator secara mekanis
\end{enumerate}

\textbf{Kenapa 500ms untuk MINOR?}\\
MINOR = 2 sensor anomali, masih ada 1 sensor normal. Artinya redundansi masih
berfungsi. Kita cukup yakin ini bukan false positive, tapi kita juga tidak
perlu terlalu lama mengunci -- 500ms cukup untuk memastikan pembacaan stabil.

\textbf{Kenapa 2000ms untuk CRITICAL?}\\
CRITICAL = 3 sensor anomali. Semua sensor ``buta''. Kita harus sangat berhati-hati
karena tidak ada informasi valid sama sekali. 2000ms = 4$\times$ lebih lama dari MINOR,
memberikan waktu lebih untuk:
\begin{itemize}[nosep]
    \item Memastikan bukan transient spike
    \item Operator manusia bisa mengambil alih jika diperlukan
    \item Proses shutdown yang aman
\end{itemize}

\textbf{Bagaimana sistem membedakan MINOR vs CRITICAL?}\\
Di simulasi, kami menggunakan durasi tekan tombol sebagai skenario:
\begin{itemize}[nosep]
    \item Klik singkat $< 5$ detik $\rightarrow$ hanya suhu dan vibrasi diinjeksi $\rightarrow$ 2/3 = MINOR
    \item Tahan $\ge 5$ detik $\rightarrow$ suhu, vibrasi, DAN tekanan diinjeksi $\rightarrow$ 3/3 = CRITICAL
\end{itemize}

Di dunia nyata, severity ditentukan oleh jumlah sensor yang melaporkan anomali
secara alami -- bukan oleh tombol. Tombol di sini adalah SIMULATOR fault injection.

\subsection{Hold-Duration Detection -- Kenapa Perlu Sticky Latch?}

\textbf{Masalah:} Di Proteus, klik tombol mouse hanya berlangsung beberapa milidetik.
Sementara polling interval kita 100ms. Tanpa mekanisme khusus, klik bisa terlewat
-- tombol HIGH hanya 10ms, keburu LOW lagi sebelum 100ms berikutnya.

\textbf{Solusi -- Sticky Latch:}\\
Begitu tombol pernah terbaca HIGH, sebuah ``latch'' (semacam flip-flop digital)
diset ke TRUE. Latch ini TETAP TRUE meskipun tombol sudah dilepas. Latch hanya
di-reset saat:
\begin{enumerate}[nosep]
    \item Fault terdeteksi (press sudah ``dikonsumsi'')
    \item Lockout selesai (sistem siap menerima press baru)
\end{enumerate}

\textbf{Kenapa tidak pakai interrupt untuk tombol?}\\
Di Proteus, interrupt tidak didukung untuk komponen push-button. Di hardware
nyata, interrupt adalah solusi yang lebih baik -- tapi untuk simulasi, sticky
latch adalah workaround yang elegan.

\textbf{Counter hold\_count:}\\
Setiap siklus di mana tombol raw HIGH, counter naik 1. Counter ini digunakan
untuk menentukan durasi tekan. Di hardware nyata, counter bisa diganti dengan
timestamp difference (waktu sekarang - waktu pertama kali ditekan).

\subsection{Event-Triggered, Power Management \& N-Scalability}

\textbf{Event-Triggered Architecture:}\\
Sistem tidak melakukan polling buta setiap interval. Logika evaluasi HANYA
dijalankan saat ada perubahan signifikan (delta) pada pembacaan sensor atau
tombol ditekan. Saat tidak ada event, sistem langsung melewati loop tanpa
evaluasi berat. Ini berbeda dengan pendekatan konvensional yang mengecek
threshold setiap siklus terlepas dari ada tidaknya perubahan.

Keuntungan:
\begin{itemize}[nosep]
    \item CPU cycles hemat -- tidak ada evaluasi yang sia-sia
    \item Latensi lebih rendah -- respons hanya saat dibutuhkan
    \item Cocok untuk sistem battery-powered
\end{itemize}

\textbf{Power Management -- Light Sleep:}\\
Saat state NORMAL dan tidak ada event, ESP32-S3 masuk ke mode light sleep.
Konsumsi daya turun drastis. Sistem dibangunkan oleh TIMG0 timer setiap
100ms untuk membaca sensor. Jika tetap NORMAL, langsung sleep kembali.

Ini penting untuk aplikasi remote/off-grid di mana power budget terbatas.

\textbf{N-Sensor Scalability:}\\
Kode didesain dengan \texttt{const N\_SENSORS: usize = 3} -- bukan hardcode ke 3 sensor.
Tinggal ganti 1 konstanta untuk skala ke 5, 7, atau 9 sensor. Voting quorum
otomatis mengikuti: $(N/2)+1$. Threshold array per sensor juga dinamis.

Artinya sistem ini bukan solusi fixed untuk 3 sensor -- ini FRAMEWORK yang
bisa diskala ke konfigurasi sensor berapa pun, tinggal sesuaikan hardware.

% ═══════════════════════════════════════════════════════════════
\section{BAGIAN 3 --- Membaca dan Memahami Data}

\subsection{Dataset Overview}

Data asli dari simulasi Proteus, 18 Juni 2026:
\begin{itemize}[nosep]
    \item \textbf{154 baris} (termasuk 8 baris header)
    \item \textbf{319 iterasi} total (iterasi 0--319, dengan heartbeat diskip)
    \item \textbf{8 fault events}: 5 MINOR, 3 CRITICAL
    \item \textbf{53 baris NORMAL + CLEARED}, \textbf{77 baris LOCKOUT}
    \item Format: \texttt{iter\ \ temp(C)\ \ press(hPa)\ \ vib(arb)\ \ latency\_us\ \ status}
\end{itemize}

\subsection{Format CSV Output}

Setiap baris output memiliki 6 kolom:

\begin{lstlisting}[style=data]
iter  temp  press  vib  latency_us  status
\end{lstlisting}

\begin{center}
\begin{tabular}{lll}
\toprule
\textbf{Kolom} & \textbf{Arti} & \textbf{Contoh} \\
\midrule
iter   & Nomor iterasi   & 4, 5, 6\dots \\
temp   & Suhu ($^{\circ}\mathrm{C}$) & 25 (normal), 99 (fault) \\
press  & Tekanan (hPa)   & 1013 (normal), 0 (CRITICAL) \\
vib    & Vibrasi         & 5 (normal), 9999 (fault) \\
latency\_us & Latensi deteksi$\rightarrow$aktuator ($\mu$s) & 45 \\
status & Status sistem   & FAULT\_DETECTED, LOCKOUT\_ACTIVE, dll \\
\bottomrule
\end{tabular}
\end{center}

\subsection{Cara Membuktikan Adaptive Lockout dari Data ASLI}

\textbf{MINOR (data asli -- iter 4--9):}

\begin{lstlisting}[style=data]
4  99 1013 9999 45 FAULT_DETECTED(MINOR,2/3)     <- fault di iter 4
5  99 1013 9999 0  LOCKOUT_ACTIVE(400ms)          <- lockout mulai
6  99 1013 9999 0  LOCKOUT_ACTIVE(300ms)
7  99 1013 9999 0  LOCKOUT_ACTIVE(200ms)
8  99 1013 9999 0  LOCKOUT_ACTIVE(100ms)
9  99 1013 9999 0  LOCKOUT_CLEARED                <- clear di iter 9
\end{lstlisting}

Dari iter 4 ke iter 9 = 5 iterasi $\times$ 100ms = \textbf{500ms}. TERBUKTI.

\textbf{CRITICAL (data asli -- iter 107--127):}

\begin{lstlisting}[style=data]
107 99 0 9999 45 FAULT_DETECTED(CRITICAL,3/3)    <- fault di iter 107
108 99 0 9999 0  LOCKOUT_ACTIVE(1900ms)           <- lockout mulai
109 99 0 9999 0  LOCKOUT_ACTIVE(1800ms)
...
126 99 0 9999 0  LOCKOUT_ACTIVE(100ms)
127 99 0 9999 0  LOCKOUT_CLEARED                  <- clear di iter 127
\end{lstlisting}

Dari iter 107 ke iter 127 = 20 iterasi $\times$ 100ms = \textbf{2000ms}. TERBUKTI.

\textbf{Perbedaan pressure -- bukti severity discrimination:}
\begin{itemize}[nosep]
    \item MINOR: press = \textbf{1013} (NORMAL -- tekanan tidak diinjeksi)
    \item CRITICAL: press = \textbf{0} (ANOMALI -- sensor tekanan gagal total)
\end{itemize}

Ini adalah BUKTI NYATA dari data mentah bahwa sistem membedakan severity
berdasarkan jumlah sensor anomali. Bukan klaim kosong -- ada di data asli.

\textbf{Konsistensi antar event -- bukti sistem deterministik:}\\
Semua 5 event MINOR menunjukkan pola yang sama persis (5 iterasi lockout).
Semua 3 event CRITICAL menunjukkan pola yang sama persis (20 iterasi lockout).
Sistem tidak pernah random -- perilakunya terprediksi dan dapat diverifikasi.

\subsection{Memahami GNUPlot}

\textbf{Panel 1 -- Multi-Sensor Readings:}\\
Garis KUNING = suhu. Garis HIJAU = vibrasi/100. Garis putus MERAH = threshold.
Spike tajam = tombol ditekan. Perhatikan: suhu selalu spike bersamaan dengan
vibrasi -- ini karena tombol menginjeksi keduanya sekaligus.

\textbf{Panel 2 -- Recovery Latency:}\\
Garis BIRU = latensi dalam mikrodetik. Semua fault punya latensi 45$\mu$s --
KONSISTEN di semua 8 event. Ini overhead MicroPython interpreter. Di Rust
bare-metal, diprediksi $< 5\mu$s karena TIMG0 punya resolusi 12.5 nanodetik.

\textbf{Panel 3 -- Status Timeline:}\\
Sumbu Y: 0=NORMAL, 1=FAULT, 2=LOCKOUT, 3=CLEARED. Pola berulang yang SAMA
setiap kali tombol ditekan = sistem deterministik. Tidak ada random behavior.

\textbf{Voting Heatmap:}\\
Tiga baris = tiga sensor. Baris atas: suhu. Tengah: tekanan. Bawah: vibrasi.
MERAH = anomali. BIRU = normal. Perhatikan: tekanan (baris tengah) HANYA merah
saat CRITICAL. Ini bukti visual bahwa CRITICAL menginjeksi 3 sensor.

\textbf{Method Comparison:}\\
Radar/spider chart membandingkan 6 dimensi. Metode kami selalu di outer ring
(nilai tertinggi) di semua dimensi. Artinya kami UNGGUL di memory safety,
concurrency, sensor fusion, fail-safe, latency, dan dokumentasi.

% ═══════════════════════════════════════════════════════════════
\section{BAGIAN 4 --- Pertanyaan Dosen \& Jawaban}

\subsection{Q1: Kenapa pakai Rust? Apa salahnya C/C++?}

\textbf{Jawaban:}\\
C/C++ tidak punya memory safety guarantee di level compiler. Programmer
bertanggung jawab penuh atas alokasi memori, pointer, dan concurrency.
Akibatnya: 186 CVE ditemukan bahkan di Rust UNSAFE code saja. Di sistem
safety-critical, bug memory bisa berakibat fatal.

Rust menjamin memory safety TANPA garbage collector -- melalui sistem ownership
yang diperiksa saat compile. Kalau ada potensi bug, kode tidak akan compile.
Ini seperti punya asisten yang memeriksa setiap baris kode sebelum dijalankan.

Rust juga zero-cost abstraction -- performanya setara C. Dibuktikan oleh
Plauska \& Liutkevi\v{c}ius (2023) yang membenchmark Rust vs C di ESP32.

\subsection{Q2: Kenapa voting $\ge 2/3$? Kenapa tidak 3/3?}

\textbf{Jawaban:}\\
Karena kita ingin TOLERANSI SATU SENSOR GAGAL. Kalau threshold 3/3, begitu
satu sensor rusak, sistem tidak akan pernah mendeteksi fault lagi (maksimal
hanya 2 sensor yang tersisa). Justru berbahaya.

Dengan threshold 2/3:
\begin{itemize}[nosep]
    \item 1 sensor rusak $\rightarrow$ masih bisa deteksi fault (asalkan 2 lainnya anomali)
    \item 2 sensor rusak $\rightarrow$ tetap deteksi fault (mayoritas)
    \item Ini adalah prinsip redundancy di safety engineering.
\end{itemize}

\subsection{Q3: Apakah sistem ini real-time?}

\textbf{Jawaban:}\\
Sistem ini SOFT real-time. Polling interval 100ms, latensi deteksi ke aktuator
45$\mu$s (di MicroPython). Di Rust bare-metal, diprediksi $< 5\mu$s.

Untuk HARD real-time, butuh interrupt-driven architecture dengan jaminan
worst-case execution time (WCET). Itu adalah arah pengembangan selanjutnya.

\subsection{Q4: Kenapa simulasi pakai MicroPython, bukan Rust langsung?}

\textbf{Jawaban:}\\
Proteus VSM mendukung MicroPython untuk ESP32-S3, tapi tidak mendukung
Rust bare-metal secara langsung. Port MicroPython adalah behavioral equivalent
-- logika voting, lockout, dan hold-duration SAMA persis dengan kode Rust.

Kode Rust sudah dicek kompilasinya (cargo check 0 error 0 warning) dan
binary-nya tersedia di GitHub. Siap di-flash ke hardware nyata.

\subsection{Q5: Dari mana nilai threshold $80^{\circ}\mathrm{C}$, 500, 900--1200?}

\textbf{Jawaban:}\\
Nilai-nilai ini adalah representative value untuk mendemonstrasikan konsep.
Di aplikasi nyata, threshold ditentukan oleh:
\begin{itemize}[nosep]
    \item Spesifikasi peralatan industri
    \item Safety analysis (HAZOP, LOPA)
    \item Regulatory requirements
\end{itemize}

Yang penting adalah MEKANISME-nya -- voting, lockout, hold-duration -- bukan
nilai threshold spesifiknya. Threshold bisa disesuaikan per aplikasi.

\subsection{Q6: Apa kontribusi orisinal riset ini?}

\textbf{Jawaban:}\\
Belum ada penelitian yang mengintegrasikan:
\begin{enumerate}[nosep,label=(\alph*)]
    \item Rust bare-metal pada ESP32-S3,
    \item Mutex<RefCell<T>> + critical\_section untuk concurrency,
    \item Voting-based 3-sensor fusion dengan majority quorum,
    \item Adaptive lockout berdasarkan fault severity,
    \item Hold-duration detection dengan sticky latch,
    \item Event-triggered architecture,
    \item N-Sensor scalability + power management,
\end{enumerate}

DALAM SATU SISTEM TERINTEGRASI.

Masing-masing komponen mungkin sudah ada di riset terpisah, tapi
KOMBINASINYA belum pernah dilakukan. Itu kontribusi kami.

\subsection{Q7: Kenapa lockout 500ms dan 2000ms? Kenapa bukan 1 detik dan 3 detik?}

\textbf{Jawaban:}\\
Nilai spesifik bisa disesuaikan. Yang penting adalah KONSEP ADAPTIVE -- durasi
lockout BERBEDA berdasarkan severity. Di dunia nyata, durasi ditentukan oleh:
\begin{itemize}[nosep]
    \item Process safety time (seberapa cepat sistem harus merespons)
    \item Valve actuation time (seberapa cepat valve bisa membuka/tutup)
    \item Process dynamics (seberapa cepat proses kembali stabil)
\end{itemize}

500ms dan 2000ms adalah nilai demonstrasi yang menunjukkan perbedaan 4$\times$ antara
MINOR dan CRITICAL. Di aplikasi nyata, nilai ini dikonfigurasi per plant.

\subsection{Q8: Apa keterbatasan sistem ini?}

\textbf{Jawaban:}\\
Jujur, ada beberapa:

\begin{enumerate}[nosep]
    \item \textbf{Overhead MicroPython:} Latensi 45$\mu$s di Proteus vs prediksi $<5\mu$s di Rust
          bare-metal. Ini bukan kelemahan desain, tapi keterbatasan platform simulasi.

    \item \textbf{Belum diuji di hardware nyata:} Semua data dari Proteus VSM. Perlu
          validasi di ESP32-S3 fisik dengan sensor sungguhan, noise lingkungan, dan
          kondisi ekstrem (EMI, suhu tinggi, power fluctuation).

    \item \textbf{Polling, bukan interrupt:} Tombol dibaca via polling sticky latch karena
          keterbatasan Proteus. Di hardware nyata, interrupt-driven akan lebih responsif
          dan hemat CPU.

    \item \textbf{Soft real-time:} Tidak ada WCET guarantee. Untuk safety-critical level
          SIL-3/4, butuh hard real-time dengan timing analysis formal.

    \item \textbf{3 sensor: masih skala kecil.} Untuk plant besar dengan puluhan sensor,
          voting quorum sederhana mungkin perlu diganti topology yang lebih kompleks.
\end{enumerate}

Yang penting: kami JUJUR tentang keterbatasan ini. Paper dan README menyebutkan
semuanya secara eksplisit. Ini riset tahap awal -- fondasinya solid, tapi
masih banyak ruang untuk dikembangkan.

% ═══════════════════════════════════════════════════════════════
\section{BAGIAN 5 --- Tips Presentasi}

\subsection{Urutan bercerita yang efektif:}

\begin{enumerate}[nosep]
    \item \textbf{BUKA dengan masalah:} ``Apa yang terjadi kalau sistem kontrol valve gagal?''
    \item \textbf{TUNJUKKAN celah riset:} ``25 jurnal, tidak ada yang mengintegrasikan
          Rust + voting + lockout''
    \item \textbf{JELASKAN solusi:} ``Safe-Concurrency Multi-Sensor Fusion -- tiga lapis
          keamanan''
    \item \textbf{DEMONSTRASIKAN bukti:} ``Ini data ASLI dari simulasi Proteus -- lihat
          perbedaan MINOR dan CRITICAL, 500ms vs 2000ms''
    \item \textbf{TUTUP dengan kontribusi + kejujuran:} ``7 elemen dalam 1 sistem. Dan kami
          tahu keterbatasannya -- ini bukan produk final, ini fondasi.''
\end{enumerate}

\subsection{Cara menjelaskan teknis ke non-teknis:}

\begin{itemize}[nosep]
    \item Ganti ``Mutex<RefCell<T>>'' $\rightarrow$ ``mekanisme penguncian otomatis''
    \item Ganti ``data race'' $\rightarrow$ ``konflik akses data bersamaan''
    \item Ganti ``critical section'' $\rightarrow$ ``zona eksklusif -- hanya satu yang boleh masuk''
    \item Ganti ``voting quorum'' $\rightarrow$ ``keputusan berdasarkan suara terbanyak''
\end{itemize}

\subsection{Cara menjelaskan data:}

\begin{itemize}[nosep]
    \item Jangan bilang ``baris 4 sampai 9''
    \item Tapi: ``Perhatikan, dari fault ke clear butuh 5 langkah. Setiap langkah 100ms.
          Jadi total 500ms. Bandingkan dengan yang ini -- 20 langkah, 2000ms.''
\end{itemize}

\subsection{Slide mana yang paling penting:}

\begin{enumerate}[nosep]
    \item Slide 2 (Celah Riset) -- buat dosen paham ini bukan proyek asal-asalan
    \item Slide 9 (Perbandingan MINOR vs CRITICAL) -- bukti utama adaptive lockout
    \item Slide 7+8 (Skenario) -- narasi ``apa yang terjadi''
    \item Slide akhir (Keterbatasan) -- tunjukin kita jujur dan paham batasan riset
\end{enumerate}

\vspace{1cm}
\begin{center}
\rule{0.5\textwidth}{0.4pt}\\[4pt]
{\small Naskah ini dibuat sebagai bahan belajar mendalam, dengan data ASLI dari simulasi}\\
{\small Proteus 18 Juni 2026. Baca, pahami, lalu jelaskan dengan kata-kata sendiri.}\\
{\small Jangan dihafal -- dimengerti.}
\end{center}

\end{document}
